(1) 利用和差化积公式，$y=2A\cos[(\Delta k\,x-\Delta\omega\,t)/2]\cos[(\bar kx-\bar\omega t)]$，其中 $\bar k=(k_1+k_2)/2$、$\bar\omega=(\omega_1+\omega_2)/2$。外层余弦控制慢变化包络，内层是快相位振荡。
(2) 一个固定包络峰满足 $\Delta k\,x-\Delta\omega\,t=2\pi n$，故 $v_g^{(2)}=dx/dt=\Delta\omega/\Delta k$。载波相位面满足 $\bar kx-\bar\omega t=$ 常数，故 $v_p=\bar\omega/\bar k$；二者一般不相同，只在色散严格线性时可能重合。
(3) 若 $\omega(k)$ 可微，由差商极限 $\lim_{\Delta k\to0}\Delta\omega/\Delta k=d\omega/dk$，所以 $v_g=\omega'(k)$，而 $v_p=\omega(k)/k$。群速度是窄带包络的几何移动速度，依赖频率对波数的局部斜率。
(4) 在色散介质中 $\omega(k)$ 并非必成正比 $k$，因而斜率与过原点的割线斜率不同，$v_g$ 可以大于或小于相速度；在有强吸收的异常色散区简单群速度甚至可超过 $c$ 或为负。这不等于可通过两列永恒单色波无损传输超光速前沿，真正信号速度要由因果响应和带宽、吸收一起判定。
